% ============ 练习问题解答 (原书 pp.133-155) ============
\chapter*{练习题解答}

\section*{第1章}

\paragraph*{练习题 1.1}
$\|a\|_{\ell^{2}} = \sqrt{\frac{n(n+1)}{2}}$, $\|a\|_{\ell^{4}} = \sqrt[4]{\frac{n(n+1)(2n+1)}{6}}$.

\paragraph*{练习题 1.2}
由 $\vec{a}+\vec{b} = (5,7,3)$, 有 $\|\vec{a}\| = \sqrt{14}$, $\|\vec{b}\| = \sqrt{29}$, $\|\vec{a}+\vec{b}\| = \sqrt{83}$ 成立. 因此, 由积分的三角不等式是显然的. $\blacksquare$

\paragraph*{练习题 1.3}
有若干种解法.

(解法 1) 利用恒等式
\[
(a^{3}+b^{3})(c^{3}+d^{3})^{2}-(ac^{2}+bd^{2})^{3} = (-bc+ad)^{2}(bc^{4}+2ac^{3}d+2bcd^{3}+ad^{4}).
\]

(解法 2) 由 Hölder 不等式, $(a^{3}+b^{3})^{\frac{1}{3}}(C^{\frac{3}{2}}+D^{\frac{3}{2}})^{\frac{2}{3}} \geq aC+bD$, $C, D > 0$ 成立, 代入 $C = c^{2}$, $D = d^{2}$, 两边三次方. $\blacksquare$

\paragraph*{练习题 1.4}
有若干种解法.

(解法 1) 画出 $s = e^{t}$ 的图形, 考察 $e^{(1-\theta)x+\theta y}$ 与 $(1-\theta)e^{x} + \theta e^{y}$ 出现在何处. 【凸性的应用】

(解法 2) 在引理~\ref{lem:1.1} 中, 取 $a = e^{(1-\theta)x}$, $b = e^{\theta y}$, $p = \frac{1}{1-\theta}$, $q = \frac{1}{\theta}$.

(解法 3) 由微积分基本定理,
\begin{align*}
(1-\theta)e^{x} + \theta e^{y} - e^{(1-\theta)x+\theta y} &= \int_{y}^{x}(1-\theta)(e^{z} - e^{(1-\theta)z+\theta y})\dd z \\
&= \int_{y}^{x}(1-\theta)(1 - e^{\theta(y-z)})e^{z}\dd z \\
&\geq 0.
\end{align*}
$\blacksquare$

\paragraph*{练习题 1.5}
对 $1 \leq p < \infty$,
\[
\|ka\|_{\ell^{p}} = \left(\sum_{n=1}^{N}|ka_{n}|^{p}\right)^{\frac{1}{p}}
\]
\[
= \left(\sum_{n=1}^{N}|k|^{p}|a_{n}|^{p}\right)^{\frac{1}{p}} = |k|\left(\sum_{n=1}^{N}|a_{n}|^{p}\right)^{\frac{1}{p}} = |k|\|a\|_{\ell^{p}}
\]
由此得到 $|k|$. $\blacksquare$

\paragraph*{练习题 1.6}
当 $a > b$ 时, $\min(a, b) = b$, $\frac{a+b-|a-b|}{2} = \frac{a+b-a+b}{2} = b$, $\max(a, b) = a$, $\frac{a+b+|a-b|}{2} = \frac{a+b+a-b}{2} = a$, 故得结论. 另一方面, 当 $a \leq b$ 时, $\min(a, b) = a$, $\frac{a+b-|a-b|}{2} = \frac{a+b-b+a}{2} = a$, $\max(a, b) = b$, $\frac{a+b+|a-b|}{2} = \frac{a+b+b-a}{2} = b$, 故得结论. $\blacksquare$

\paragraph*{练习题 1.7}
(1) $\infty$, (2) 1, (3) $\infty$, (4) $\infty$, (5) $\infty$, (6) 2.

\paragraph*{练习题 1.8}
图形从略, 积分值为 $-3$. $\blacksquare$

\paragraph*{练习题 1.9}
(1) 由定理~\ref{thm:1.3},
\[
\sum_{n=1}^{N}|a_{n}| = \int_{0}^{\infty}\sharp\{1 \leq n \leq N: |a_{n}| > t\}\dd t
\]
成立. 这里, 当 $t > 0$ 时, $|a_{n}| > t$ 与 $|a_{n}| > [t]$ 等价, 因此
\[
\sum_{n=1}^{N}|a_{n}| = \int_{0}^{\infty}\sharp\{1 \leq n \leq N: |a_{n}| > [t]\}\dd t.
\]
将积分分割来考虑,
\begin{align*}
\int_{0}^{\infty}\sharp\{1 \leq n \leq N: |a_{n}| > [t]\}\dd t &= \sum_{k=0}^{\infty}\int_{k}^{k+1}\sharp\{1 \leq n \leq N: |a_{n}| > k\}\dd t \\
&= \sum_{k=0}^{\infty}\sharp\{1 \leq n \leq N: |a_{n}| > k\}.
\end{align*}

(2) 设 $a_{n}$ 为将 $n$ 素因数分解所得因子 $5$ 的个数, 则 $n!$ 用十进制书写时末尾出现的 $0$ 的个数表为
\[
\sum_{k=1}^{n}a_{k}.
\]
由 (1),
\[
\sum_{k=1}^{n}a_{k} = \sum_{k=0}^{\infty}\sharp\{1 \leq m \leq n: a_{m} > k\}
\]
成立. 这里, 由于 $a_{n} > k$ 与 $a_{n} \geq k+1$ 等价, 且与 $n$ 含有因子 $5^{k+1}$ 等价, 所以
\[
\sharp\{1 \leq m \leq n: a_{m} > k\} = \left[\frac{n}{5^{k+1}}\right].
\]
将其代入, 即得结论. $\blacksquare$

\paragraph*{练习题 1.10}
(1) 1.

(2) 2.

(3) 算子范数由 $\sup\{\sqrt{3x^{2}+2xy+y^{2}} : x^{2}+y^{2} = 1\}$ 给出. 用三角函数计算, 该值为 $\sup\{\sqrt{2+\sqrt{2}\cos(\theta+\pi/4)} : \theta \in \mathbb{R}\} = \sqrt{2+\sqrt{2}}$. $\blacksquare$

\paragraph*{练习题 1.11}
(1) 当 $x^{2}+y^{2} = 1$ 时, $(4x-2y)^{2}+(8x-4y)^{2} = 20(2x-y)^{2} \leq 100$, 故算子范数为 $10$.

(2) 因 $f^{2} = 0$, 算子范数为 $0$.

(3) $\|f^{2}\| = 0 < 100 = \|f\|^{2}$.

\paragraph*{练习题 1.12}
给出相应的条件最大值问题.

(1) 求当 $\sqrt{x^{2}+y^{2}} = 1$ 时 $\sqrt{(4x-3y)^{2}+(x-4y)^{2}}$ 的最大值. 其最大值为 $\sqrt{21+4\sqrt{17}} = 2+\sqrt{17}$.

(2) 求当 $|x|+|y| = 1$ 时 $|4x-3y|+|x-4y|$ 的最大值. 当 $x = 0$, $y = 1$ 时为 7. 当 $|x|+|y| = 1$ 时, 由 $|4x-3y|+|x-4y| \leq 5|x|+7|y| \leq 7(|x|+|y|) = 7$, 为 7.

(3) 求当 $\sqrt{x^{2}+y^{2}} = 1$ 时 $\max(|4x-3y|, |x-4y|)$ 的最大值. 当 $x^{2}+y^{2} = 1$ 时, $|4x-3y|$ 的最大值为 $5$. 同样, $|x-4y|$ 的最大值为 $\sqrt{17}$. 因此, 算子范数为 $5$.

(4) 求当 $\max(|x|, |y|) = 1$ 时 $\sqrt{(4x-3y)^{2}+(x-4y)^{2}}$ 的最大值. 实际上, 因 $(4x-3y)^{2} \leq 49$, $(x-4y)^{2} \leq 25$, 所求值为 $\sqrt{74}$. 范数值在 $(x,y) = (1,-1)$ 处取得. $\blacksquare$

\paragraph*{练习题 1.13}
在条件 $x^{2}+y^{2} = 1$ 下求 $(2x-3y)^{2}+(3x+4y)^{2} = 13x^{2}+12xy+25y^{2}$ 的最大值. 若 $x = \cos\theta$, $y = \sin\theta$, 则 $13x^{2}+12xy+25y^{2} = 19-6\cos 2\theta+6\sin 2\theta$, 故第一个算子范数为 $\sqrt{19+6\sqrt{2}} = 1+3\sqrt{2}$. 又, 在条件 $\max(|x|, |y|) = 1$ 下 $(2x-3y)^{2}+(3x+4y)^{2} = 13x^{2}+12xy+25y^{2}$ 的最大值为 50, 故第二个算子范数为 $5\sqrt{2}$. $\blacksquare$

\paragraph*{练习题 1.14}
在条件 $x^{2}+y^{2} = 1$ 下, $(2x+y)^{2}+(x-2y)^{2}+(x+y)^{2} = 6x^{2}+6y^{2}+2xy$ 的最大值为 7, 故所求算子范数为 $\sqrt{7}$. $\blacksquare$

\paragraph*{练习题 1.15}
由 Hölder 不等式, $\|(2,1)\|_{\ell^{3/2}} = (2\sqrt{2}+1)^{\frac{2}{3}}$. $\blacksquare$

\paragraph*{练习题 1.16}
设 $\vec{p} = (x,y)$, $\vec{a} = (1,-4)$, $\vec{b} = (2,-3)$, 由 Minkowski 不等式,
\[
\|\vec{p}+\vec{a}\|_{\ell^{3}} \leq \|\vec{p}+\vec{b}\|_{\ell^{3}} + \|\vec{a}-\vec{b}\|_{\ell^{3}}
\]
成立. 因此 $f(x,y) \leq \|\vec{a}-\vec{b}\|_{\ell^{3}} = \sqrt[3]{2}$. 这里, 若 $\vec{p} = -\vec{b}$, 则等号成立, 故最大值为 $\sqrt[3]{2}$. $\blacksquare$

\paragraph*{练习题 1.17}
(1) 对实数列 $a = \{a_{n}\}_{n=1}^{N}$,
\[
\|U_{N}a\|_{\ell^{1}} = \sum_{n=1}^{N}|b_{n}| = \sum_{n=1}^{N}\frac{|a_{1}+a_{2}+\cdots+a_{n}|}{n} \leq \sum_{n=1}^{N}\frac{\|a\|_{\ell^{1}}}{n}
\]
由 $(1,0,0,\ldots,0)$ 达成等号, 故 $\|U_{N}\|_{\ell^{1}\to\ell^{1}} = \sum_{n=1}^{N}\frac{1}{n}$.

(2) 取平均的操作不增大最大值, 故显然.

(3) 当 $\lambda \geq 1$ 或 $\lambda \leq \frac{1}{N}$ 时, 由所考虑的集合分别为空集与 $\{1,2,\ldots,N\}$ 的子集可知是显然的. 设 $\frac{1}{N} < \lambda < 1$. 取满足 $\frac{1}{l} < \lambda \leq \frac{1}{l-1}$ 的自然数 $l = 2, 3, \ldots, N$. 于是,
\[
\{k \in \{1, 2, \ldots, N\}: |b_{k}| > \lambda\|a\|_{\ell^{1}}\} \subset \left\{k \in \{1, 2, \ldots, N\}: |b_{k}| > \frac{1}{l}\|a\|_{\ell^{1}}\right\} \subset \{1, 2, \ldots, l-1\}
\]
所以
\[
\sharp\{k \in \{1, 2, \ldots, N\}: |b_{k}| > \lambda\|a\|_{\ell^{1}}\} \leq l-1 \leq \frac{1}{\lambda}.
\]

(4) 用 Marcinkiewicz 插值定理对 (2) 与 (3) 进行插值.

(5) 设字母变量 $a$, $b$, $c$ 为实数. 设 $S = |a|^{2}+|b|^{2}+|c|^{2}$. 由
\begin{align*}
\|U_{1}\|_{\ell^{2}\to\ell^{2}} &= \max\{|a|: |a| = 1\} = 1, \\
\|U_{2}\|_{\ell^{2}\to\ell^{2}} &= \max\{\sqrt{|a|^{2}+4^{-1}|a+b|^{2}}: |a|^{2}+|b|^{2} = 1\} \\
&= \frac{1+\sqrt{5}}{2\sqrt{2}} < 1.15,
\end{align*}
\begin{align*}
\|U_{3}\|_{\ell^{2}\to\ell^{2}} &= \max\{\sqrt{|a|^{2}+4^{-1}|a+b|^{2}+9^{-1}|a+b+c|^{2}}: S = 1\} \\
&\geq \|U_{3}\mathbf{e}_{1}\|_{\ell^{2}} = \frac{7}{6} > 1.15.
\end{align*}

\paragraph*{练习题 1.18}
由
\[
r_{1}(t)r_{2}(t)r_{3}(t) = \chi_{[0,1/8)\cup[3/8,1/2)\cup[5/8,7/8)}(t) - \chi_{[1/8,3/8)\cup[1/2,5/8)\cup[7/8,1)}(t),
\]
所求积分值为 $-\frac{1}{8}$. $\blacksquare$

\paragraph*{练习题 1.19}
由
\[
2e^{t^{2}} = \sum_{n=0}^{\infty}\frac{2}{n!}t^{2n}, \quad e^{t}+e^{-t} = \sum_{n=0}^{\infty}\frac{2}{(2n)!}t^{2n}
\]
是显然的. $\blacksquare$

\paragraph*{练习题 1.20}
只需在 $0 \leq t < 1$ 的范围内考虑, 所求为 $\sum_{j=1}^{N}|a_{j}|$. $\blacksquare$

\paragraph*{练习题 1.21}
\begin{align*}
\mathcal{U}_{0} &= \{\{1\}, \{3\}, \{6\}, \{8\}, \{12\}, \{13\}, \{8,9,10,11,12,13,14,15\}, \\
&\qquad \{14\}, \{6,7\}, \{8,9\}, \{12,13\}, \{14,15\}, \{8,9,10,11\}, \{12,13,14,15\}, \\
&\qquad \{0,1,2,3,4,5,6,7,8,9,10,11,12,13,14,15\}\},
\end{align*}
\[
\mathcal{U} = \{\{0,1,2,3,4,5,6,7,8,9,10,11,12,13,14,15\}\}, \quad E = I(4,0),
\]
\[
g_{0} = g_{1} = \cdots = g_{15} = 3,
\]
\[
b: -3, 1, -3, 1, -1, -3, 3, -3, 9, -3, -1, -3, 3, 3, 3, -3,
\]
\[
b_{\{0,1,2,3,4,5,6,7,8,9,10,11,12,13,14,15\}} = b.
\]
$\blacksquare$

\paragraph*{练习题 1.22}
\begin{itemize}
\item $E_{0}(a)$: $a$ 本身
\item $E_{2}(a)$: $2, 2, 2, 2, 2, 2, 2, 2, 3.5, 3.5, 3.5, 3.5, 4.5, 4.5, 4.5, 4.5$
\item $E_{3}(a)$: $2, 2, 2, 2, 2, 2, 2, 2, 4, 4, 4, 4, 4, 4, 4, 4$
\item $E_{4}(a)$: $3, 3, 3, 3, 3, 3, 3, 3, 3, 3, 3, 3, 3, 3, 3, 3$
\item $D_{0}(a)$: $-2, 2, -2, 2, 1, -1, 3, -3, 6, -6, 1, -1, 0, 0, 3, -3$
\item $D_{1}(a)$: $0, 0, 0, 0, -1, -1, 1, 1, 2.5, 2.5, -2.5, -2.5, 1.5, 1.5, -1.5, -1.5$
\item $D_{2}(a)$: $0, 0, 0, 0, 0, 0, 0, 0, -0.5, -0.5, -0.5, -0.5, 0.5, 0.5, 0.5, 0.5$
\item $D_{3}(a)$: $-1, -1, -1, -1, -1, -1, -1, -1, 1, 1, 1, 1, 1, 1, 1, 1$ $\blacksquare$
\end{itemize}

\paragraph*{练习题 1.23}
如直接计算所见, 为 $\sum_{k=0}^{7}a_{k}^{2}$, 但这可视为定理~\ref{thm:1.9}(3) 中 $N = 3$ 情形的结果. 实际上,
\[
E_{3}(a) = \{A, A, A, A, A, A, A, A\},
\]
\[
D_{0}(a) = \{B, -B, C, -C, D, -D, E, -E\},
\]
\[
D_{1}(a) = \{F, F, -F, -F, G, G, -G, -G\},
\]
\[
D_{2}(a) = \{H, H, H, H, -H, -H, -H, -H\}
\]
故 $\|a\|_{\ell^{2}}^{2} = \|E_{N}(a)\|_{\ell^{2}}^{2}+\sum_{k=0}^{N-1}\|D_{k}(a)\|_{\ell^{2}}^{2}$ 的右边为 $8A^{2}+2B^{2}+2C^{2}+2D^{2}+2E^{2}+4F^{2}+4G^{2}+8H^{2}$, 左边为 $\sum_{k=0}^{7}a_{k}^{2}$. $\blacksquare$

\paragraph*{练习题 1.24}
对常数数列 $m^{k} = (-1)^{k}$ 使用定理~\ref{thm:1.12}. $\blacksquare$

\paragraph*{练习题 1.25}
\begin{align*}
\|g\|_{\ell^{2}} &= \sqrt{\sum_{j=1}^{N}|g_{j}|^{2}} \leq \sqrt{\sup_{j=1,2,\ldots,N}|g_{j}|\sum_{j=1}^{N}|g_{j}|} \\
&= \sqrt{\|g\|_{\ell^{1}}\|g\|_{\ell^{\infty}}}.
\end{align*}

\paragraph*{练习题 1.26}
关于前半部分, 由 $E_{k}(e)_{j} = 2^{-k}$, $E_{k+1}(e)_{j} = 2^{-k-1}$ 是显然的. 后半部分如下证明.
\begin{align*}
\|g(e)\|_{\ell^{1}} &= \sum_{j=0}^{2^{N}-1}|g(e)_{j}| \\
&\geq |E_{N}(e)| + \sum_{k=0}^{N-1}\sum_{j=2^{k}}^{2^{k+1}-1}|D_{k}(e)_{j}| \\
&\geq \sum_{k=0}^{N-1}\sum_{j=2^{k}}^{2^{k+1}-1}2^{-k-2} \\
&= \frac{N}{4}.
\end{align*}

\paragraph*{练习题 1.27}
设 $a = \{a_{j_{1}j_{2}}\}_{j_{1},j_{2}=0,1,\ldots,2^{N}-1}$, $E_{k}^{2}(a) = \{b_{j_{1}j_{2}}\}_{j_{1},j_{2}=0,1,\ldots,2^{N}-1}$, 则
\[
b_{j_{1}j_{2}} = \frac{1}{2^{k}}\sum_{\mu=2^{k}[2^{-k}l]}^{2^{k}[2^{-k}l]+2^{k}-1}a_{j_{1}\mu}.
\]
又, 设 $E_{k'}^{1}(E_{k}^{2}(a)) = \{c_{j_{1}j_{2}}\}_{j_{1},j_{2}=0,1,\ldots,2^{N}-1}$, 则
\begin{align*}
c_{j_{1}j_{2}} &= \frac{1}{2^{k'}}\sum_{\nu=2^{k'}[2^{-k'}l]}^{2^{k'}[2^{-k'}l]+2^{k'}-1}a_{\nu j_{2}} \\
&= \frac{1}{2^{k'+k}}\sum_{\mu=2^{k}[2^{-k}l]}^{2^{k}[2^{-k}l]+2^{k}-1}\sum_{\nu=2^{k'}[2^{-k'}l]}^{2^{k'}[2^{-k'}l]+2^{k'}-1}a_{\nu\mu}.
\end{align*}
因此, 可知关于 $\mu$, $\nu$ 的求和顺序可以交换, 故得 $E_{k}^{2} \circ E_{k'}^{1}(a) = E_{k'}^{1} \circ E_{k}^{2}(a)$. $\blacksquare$

\paragraph*{练习题 1.28}
固定 $j$, 任意取 $J < j$. 由
\[
ja_{j} \leq Ja_{j} + \sum_{k=J+1}^{j}a_{k}
\]
令 $j \to \infty$, 得
\[
\limsup_{j \to \infty}ja_{j} \leq \sum_{k=J+1}^{\infty}a_{k}.
\]
令 $J \to \infty$ 即得结论. $\blacksquare$

\paragraph*{练习题 1.29}
(1) $2^{\frac{1}{p}}$.

(2) 若其为紧集, 则本应能从 $\{\mathbf{e}_{j}\}_{j=0}^{\infty}$ 中取出收敛子列, 但由 (1), 任何子列都不收敛, 故矛盾. $\blacksquare$

\paragraph*{练习题 1.30}
(1) 由 $x_{k} - x_{k'}$ 的具体形状
\[
x_{k} - x_{k'} = (0, 0, \ldots, 0, 2^{-k-1}, 2^{-k-2}, \ldots, 2^{-k'}, 0, 0, \ldots)
\]
是显然的.

(2) 对任意 $x = (x^{0}, x^{1}, \ldots) \in \operatorname{Span}(\{\mathbf{e}_{j}\}_{j=0}^{\infty})$, 设 $k = \max\{j \in \mathbb{N}_{0} : x^{j} \neq 0\}$. 当 $j > k$ 时, $\|x - x_{k}\|_{\ell^{\infty}} \geq 2^{-k}$, 故 $\{x_{k}\}_{k=1}^{\infty}$ 不收敛. $\blacksquare$

\section*{第2章}

\paragraph*{练习题 2.1}
设 $A \leq B$. 此时,
\[
2^{-1}(A^{p}+B^{p}) \leq B^{p} \leq (A+B)^{p} \leq (2B)^{p} \leq 2^{p}(A^{p}+B^{p})
\]
故得结论. $\blacksquare$

\paragraph*{练习题 2.2}
(1) $\|f\|_{L^{1}(0,\infty)} = \int_{0}^{\infty}\frac{\dd x}{1+x} = [\log(1+x)]_{0}^{\infty} = \infty$.

(2) $\|f\|_{L^{1}(0,\infty)} = \int_{0}^{\infty}\frac{\dd x}{(1+x)^{a}} = \left[\frac{(1+x)^{1-a}}{1-a}\right]_{0}^{\infty} = \infty$.

(3) $\|f\|_{L^{1}(0,\infty)} = \int_{0}^{\infty}\frac{\dd x}{(1+x)^{a}} = \left[\frac{(1+x)^{1-a}}{1-a}\right]_{0}^{\infty} = \frac{1}{a-1}$.

(4) 由 (1), (2), (3), $\|f\|_{L^{1}(0,\infty)} < \infty$ 成立的必要充分条件是 $a > 1$.

(5) 由 (4), $2a > 1$, 即 $a > \frac{1}{2}$. $\blacksquare$

\paragraph*{练习题 2.3}
省略 $L^{p}(\mathbb{R}^{n})$ 的定义. 为证明 $f \in L^{p}(\mathbb{R}^{n})$, 需要证明
\[
\int_{\mathbb{R}^{n}}\chi_{\mathbb{R}^{n}\setminus B(1)}(x)|x|^{ap}\dd x < \infty,
\]
即
\[
\int_{\mathbb{R}^{n}\setminus B(1)}|x|^{ap}\dd x < \infty.
\]
由 $\mathbb{R}^{n}\setminus B(1) = \bigcup_{j=1}^{\infty}B(2^{j})\setminus B(2^{j-1})$ 与变量替换 $x \mapsto 2^{-j}x$,
\begin{align*}
\int_{\mathbb{R}^{n}\setminus B(1)}|x|^{ap}\dd x &= \sum_{j=1}^{\infty}\int_{B(2^{j})\setminus B(2^{j-1})}|x|^{ap}\dd x \\
&= \sum_{j=1}^{\infty}2^{j(ap+n)}\int_{B(1)\setminus B(2^{-1})}|x|^{ap}\dd x \\
&= \frac{2^{ap+n}}{1-2^{ap+n}}\int_{B(1)\setminus B(2^{-1})}|x|^{ap}\dd x < \infty
\end{align*}
即得. 从该级数的计算可以看出, 若换为 $a = -n/p$, 则最后一个式子的分母变为 $0$, 论证便行不通了. $\blacksquare$

\section*{第3章}

\paragraph*{练习题 3.1}
(1) 画图便知, $[-1,1]^{n} \cap Q(x, |x|+2) = [-1,1]^{n}$. 因此,
\begin{align*}
M\chi_{[-1,1]^{n}}(x) &= \sup_{Q: Q \text{ 是立方体}}\frac{|Q \cap [-1,1]^{n}|}{|Q|} \\
&\geq \frac{|[-1,1]^{n} \cap Q(x, |x|+2)|}{|Q(x, |x|+2)|} \\
&= \frac{2^{n}}{(|x|+2)^{n}}.
\end{align*}

(2) 由提示中的不等式,
\begin{align*}
\int_{\mathbb{R}^{n}}Mf(x)\dd x &\geq \int_{\mathbb{R}^{n}}\frac{2^{n}}{(|x|+2)^{n}}\dd x \\
&\geq \int_{\mathbb{R}^{n}}\prod_{j=1}^{n}\frac{2}{(|x_{j}|+2)}\dd x \\
&= \prod_{j=1}^{n}\int_{\mathbb{R}}\frac{2}{(|x_{j}|+2)}\dd x_{j} \\
&= \infty
\end{align*}
即得. $\blacksquare$

\paragraph*{练习题 3.2}
如提示所述用 $A$ 表示条件等, 则要判定的条件与
\[
\varepsilon < \sqrt[n]{\frac{A}{\lambda}} - 1
\]
相同. 因此, 可知可取
\[
\varepsilon = \min\left(1, \frac{1}{2}\left(\sqrt[n]{\frac{A}{\lambda}} - 1\right)\right).
\]
$\blacksquare$

\paragraph*{练习题 3.3}
从略.

\paragraph*{练习题 3.4}
因 $B$ 最长, 取 $K_{1} = B$. 在与 $B$ 不相交者中最长的是 $E$, 故取 $K_{2} = E$. 因 $A \subset 3B$, $C \subset 3B$, $D \subset 3B$, $F \subset 3B$, 取 $L = 2$, $K_{1} = B$, $K_{2} = E$, $\iota(A) = \iota(B) = \iota(C) = \iota(D) = \iota(F) = 1$, $\iota(E) = 2$ 即可. 【注意】也有取 $K_{1} = E$ 的方法. $\blacksquare$

\paragraph*{练习题 3.5}
(1) 由
\[
b_{n,k}(x) - b_{n,k+1}(x) = \frac{n!}{(k+1)!(n-k)!}x^{k}(1-x)^{n-k-1}(k+1-nx-x)
\]
是显然的.

(2) 由 (1), 在 $t \leq \frac{k_{n,x}}{n+1}$ 上 $K(x,\cdot)$ 单调递增, 在 $\frac{k_{n,x}}{n+1} \leq t \leq \frac{k_{n,x}+1}{n+1}$ 上为常数函数, 在 $t \geq \frac{k_{n,x}+1}{n+1}$ 上单调递减.

(3) 计算右边,
\begin{align*}
\int_{[0,1]}F(t)\dd\mu_{x}(t) &= (n+1)\sum_{k=0}^{n}b_{n,k}(x)\left\{F\left(\frac{k+1}{n+1}\right) - F\left(\frac{k}{n+1}\right)\right\} \\
&= (n+1)\sum_{k=0}^{n}b_{n,k}(x)\int_{\frac{k}{n+1}}^{\frac{k+1}{n+1}}f(t)\dd t \\
&= K_{n}f(x, t).
\end{align*}

(4) 因 $|K_{n}f(x,t)| \leq K_{n}[|f|](x,t)$, 若必要, 可考虑以 $|f|$ 代替 $f$, 设 $f \geq 0$. 对 (3) 中的 $F$, 若 $t \geq x$, 则 $0 \leq F(t) \leq (t-x)Mf(x)$; 若 $t \leq x$, 则 $0 \geq F(t) \geq (t-x)Mf(x)$. 又由 (1), 当 $t \geq x$ 时 $\mu_{x} \geq 0$, 当 $t \leq x$ 时 $\mu_{x} \leq 0$. 因此由 (3), $0 \leq K_{n}f(x,t) \leq \int_{[0,1]}(t-x)Mf(x)\dd\mu_{x}(t)$. 将 (3) 应用于常数函数 $Mf(x)$ 以代替 $f$, 得 $0 \leq K_{n}f(x,t) \leq K_{n}[Mf(x)](x,t) = Mf(x)$. $\blacksquare$

\paragraph*{练习题 3.6}
(1) 任意取包含 $x$ 的立方体 $Q$. 由积分的三角不等式,
\[
\frac{1}{|Q|}\int_{Q}|f(y)+g(y)|\dd y \leq \frac{1}{|Q|}\int_{Q}|f(y)|\dd y + \frac{1}{|Q|}\int_{Q}|g(y)|\dd y.
\]
由 $Mf(x)$, $Mg(x)$ 的定义,
\[
\frac{1}{|Q|}\int_{Q}|f(y)|\dd y + \frac{1}{|Q|}\int_{Q}|g(y)|\dd y \leq Mf(x) + Mg(x).
\]
将这些组合起来, 得
\[
\frac{1}{|Q|}\int_{Q}|f(y)+g(y)|\dd y \leq Mf(x) + Mg(x).
\]
对包含 $x$ 的立方体 $Q$ 取上确界, 即得结论.

(2) 由绝对值的性质,
\[
\frac{1}{|Q|}\int_{Q}|af(y)|\dd y = |a| \times \frac{1}{|Q|}\int_{Q}|f(y)|\dd y
\]
故结论显然. $\blacksquare$

\paragraph*{练习题 3.7}
(1) \textbf{【当 $a < 5$ 时】}
\begin{align*}
F_{a}(t) &= \int_{5}^{\min(8, \max(5, t))}(s-5)\dd s \\
&= \frac{1}{2}(\min(8, \max(5, t))-5)^{2} \\
&= \frac{1}{2}\min(3, \max(0, t-5))^{2}.
\end{align*}
若 $t < 5$, 则 $F_{a}(t) = 0$, 故在 $t = a$ 处可微. 以后, 设
\[
\Psi(t) = \frac{1}{2}\min(3, \max(0, t-5))^{2}.
\]
此时, 对一切实数 $t$, $F_{a}(t) = \Psi(t) - \Psi(a)$.

\textbf{【当 $a = 5$ 时】} $F_{5}(t) = \Psi(t)$, 由 $\Psi'(t-0) = 0$, $\Psi'(t+0) = 0$ 可微.

\textbf{【当 $5 < a < 8$ 时】} 当 $t \in (5, 8)$ 时, $\Psi(t) = \frac{1}{2}(t-5)^{2}$, 故 $\Psi$ 在 $t = a$ 处可微. 因此 $F_{a}$ 也在 $t = a$ 处可微.

\textbf{【当 $a = 8$ 时】} 若 $t < 8$, 则 $F_{8}(t) = \frac{1}{2}(t-5)^{2} - \frac{9}{2}$. 因此, 若 $t < 8$, 则 $F_{8}'(8-0) = 3$. 另一方面, 若 $t > 8$, 则 $F_{8}(t) = -\frac{9}{2}$, 故 $F_{8}'(8+0) = 0$. 右导数与左导数不一致, 故 $F_{8}$ 在 $t = 8$ 处不可微.

\textbf{【当 $a > 8$ 时】} 当 $t > 8$ 时, $F_{a}(t) = \Psi(t) - \Psi(a) = -\frac{9}{2} + \frac{9}{2} = 0$, 故 $F_{a}'(a) = 0$.

(2) \textbf{【当 $a < 6$ 时】}
\begin{align*}
F_{a}(t) &= \int_{6}^{\min(9, \max(6, t))}(s-9)\dd s \\
&= \frac{1}{2}(\min(9, \max(6, t))-9)^{2} - \frac{9}{2} \\
&= \frac{1}{2}\min(0, \max(-3, t-9))^{2} - \frac{9}{2}.
\end{align*}
若 $t < 6$, 则 $F_{a}(t) = 0$, 故在 $t = a$ 处可微. 以后, 设
\[
\Psi(t) = \frac{1}{2}\min(0, \max(-3, t-9))^{2} - \frac{9}{2}.
\]
此时, 对一切实数 $t$, $F_{a}(t) = \Psi(t) - \Psi(a)$.

\textbf{【当 $a = 6$ 时】} 若 $t < 6$, 则 $F_{6}(t) = 0$. 因此, 若 $t < 6$, 则 $F_{6}'(6-0) = 0$. 另一方面, 若 $t > 6$, 则 $F_{6}(t) = \frac{1}{2}(t-9)^{2} - \frac{9}{2}$, 故 $F_{6}'(6+0) = -3$. 右导数与左导数不一致, 故 $F_{6}$ 在 $t = 6$ 处不可微.

\textbf{【当 $6 < a < 9$ 时】} 当 $t \in (6, 9)$ 时, $\Psi(t) = \frac{1}{2}(t-9)^{2}$, 故 $\Psi$ 在 $t = a$ 处可微. 因此 $F_{a}$ 也在 $t = a$ 处可微.

\textbf{【当 $a = 9$ 时】} $F_{9}(t) = \Psi(t) - \Psi(9)$, 由 $F_{9}'(t-0) = 0$, $F_{9}'(t+0) = 0$ 可微.

\textbf{【当 $a > 9$ 时】} 当 $t > 9$ 时, $F_{a}(t) = \Psi(t) - \Psi(a) = -\frac{9}{2} + \frac{9}{2} = 0$, 故 $F_{a}'(t) = 0$. $\blacksquare$

\paragraph*{练习题 3.8}
(1) 由 $\frac{1}{r}\int_{x}^{x+r}|f(y)-f(x)|\dd y \leq \frac{1}{r}\int_{x-r}^{x+r}|f(y)-f(x)|\dd y$ 与定理~\ref{thm:3.9}, 结论显然.

(2) 用 $v_{n}$ 表示 $n$ 维单位球 $B(1)$ 的体积. 由 $Q(x,r) \supset B(x,r)$, $|Q(x,r)| : |B(x,r)| = 2^{n} : v_{n}$,
\[
\frac{1}{|B(x,r)|}\int_{B(x,r)}|f(y)-f(x)|\dd y \leq \frac{2^{n}}{v_{n}|Q(x,r)|}\int_{Q(x,r)}|f(y)-f(x)|\dd y
\]
再次由定理~\ref{thm:3.9}, 结论显然. $\blacksquare$

\paragraph*{练习题 3.9}
(1) 注意 $Q(x,r) \supset B(x,r)$. 利用不等式
\[
\frac{1}{|B(x,r)|}\int_{B(x,r)}|f(y)-f(x)|\dd y \leq \frac{|Q(0,1)|}{|B(0,1)|} \times \frac{1}{|Q(x,r)|}\int_{Q(x,r)}|f(y)-f(x)|\dd y
\]
得
\begin{align*}
\lim_{r \downarrow 0}\frac{1}{|B(x,r)|}\int_{B(x,r)}|f(y)-f(x)|\dd y &\leq \frac{|Q(0,1)|}{|B(0,1)|}\lim_{r \downarrow 0}\frac{1}{|Q(x,r)|}\int_{Q(x,r)}|f(y)-f(x)|\dd y \\
&= 0.
\end{align*}

(2) 因 $Q(x,r) \subset B(x, \sqrt{n}r)$,
\begin{align*}
\frac{1}{|Q(x,r)|}\int_{Q(x,r)}|f(y)-f(x)|\dd y &\leq \frac{1}{|Q(x,r)|}\int_{B(x, \sqrt{n}r)}|f(y)-f(x)|\dd y \\
&= \frac{|B(0, \sqrt{n})|}{|Q(0,1)|} \times \frac{1}{|B(x, \sqrt{n}r)|}\int_{B(x, \sqrt{n}r)}|f(y)-f(x)|\dd y
\end{align*}
即得. 因此,
\[
\lim_{r \downarrow 0}\frac{1}{|Q(x,r)|}\int_{Q(x,r)}|f(y)-f(x)|\dd y \leq \frac{|B(0, \sqrt{n})|}{|Q(0,1)|}\lim_{r \downarrow 0}\frac{1}{|B(x, \sqrt{n}r)|}\int_{B(x, \sqrt{n}r)}|f(y)-f(x)|\dd y = 0.
\]
$\blacksquare$

\paragraph*{练习题 3.10}
若这样的 $E$ 存在, 则由 $E$ 所满足的条件, $\lim_{r\downarrow 0}\frac{|E\cap(x-r,x+r)|}{2r} = \frac{1}{2}$. 另一方面, 由 Lebesgue 微分定理, 在 $\chi_{E}$ 的 Lebesgue 点处应有 $\lim_{r\downarrow 0}\frac{|E\cap(x-r,x+r)|}{2r} = 0, 1$, 故矛盾. $\blacksquare$

\section*{第4章}

\paragraph*{练习题 4.1}
对 $m$ 用数学归纳法证明. 等式在 $m = 1$ 的情形就是定理~\ref{thm:2.31}(2) 本身. 由定理~\ref{thm:2.31}(1), $A^{*}A$ 是自伴的. 设 $M \in \mathbb{N}$, 若 $(A^{*}A)^{2^{M}}$ 自伴且 $\|(A^{*}A)^{2^{M}}\| = \|A\|^{2^{M+1}}$ 成立, 则 $(A^{*}A)^{2^{M}}$ 自伴, 由定理~\ref{thm:2.31}(2),
\begin{align*}
\|(A^{*}A)^{2^{M+1}}\| &= \|(A^{*}A)^{2^{M}}(A^{*}A)^{2^{M}}\| \\
&= \|(A^{*}A)^{2^{M}}\| \cdot \|(A^{*}A)^{2^{M}}\| \\
&= (\|A\|^{2^{M}})^{2} \\
&= \|A\|^{2^{M+1}}
\end{align*}
成立. 因此, 由数学归纳法, 证明完成. $\blacksquare$

\paragraph*{练习题 4.2}
由 $f \in C_{\mathrm{c}}^{\infty}(\mathbb{R}^{n})$, 积分
\[
\int_{\mathbb{R}^{n}\setminus B(x,1)}\frac{(x_{j}-y_{j})f(y)}{|x-y|^{n+1}}\dd y
\]
绝对收敛. 因此, 只需证明
\[
\lim_{\varepsilon \downarrow 0}\int_{B(x,1)\setminus B(x,\varepsilon)}\frac{(x_{j}-y_{j})f(y)}{|x-y|^{n+1}}\dd y
\]
存在. 如提示所述,
\[
\int_{B(x,1)\setminus B(x,\varepsilon)}\frac{(x_{j}-y_{j})f(x)}{|x-y|^{n+1}}\dd y = 0
\]
所以只需证明
\[
\lim_{\varepsilon \downarrow 0}\int_{B(x,1)\setminus B(x,\varepsilon)}\frac{(x_{j}-y_{j})(f(y)-f(x))}{|x-y|^{n+1}}\dd y
\]
存在. 这里由中值定理,
\begin{align*}
\int_{B(x,1)\setminus B(x,\varepsilon)}\frac{|(x_{j}-y_{j})(f(y)-f(x))|}{|x-y|^{n+1}}\dd y &\leq \sup_{z \in \mathbb{R}^{n}}|\nabla f(z)|\int_{B(x,1)\setminus B(x,\varepsilon)}\frac{1}{|x-y|^{n-1}}\dd y \\
&\leq \sup_{z \in \mathbb{R}^{n}}|\nabla f(z)|\int_{B(x,1)}\frac{1}{|x-y|^{n-1}}\dd y < \infty
\end{align*}
故由 Lebesgue 收敛定理极限存在. $\blacksquare$

\paragraph*{练习题 4.3}
(1) 因 $\varphi'(t) = 3(2-t)\max(0,2t-t^{2})^{2}$, $\varphi'$ 是 $C^{1}$ 级的.

(2) 与 (3) 实际上, 令 $\rho(0) = 0$, 设
\[
\rho(t) = \varphi(t)\left(\sum_{j=-\infty}^{\infty}\varphi(2^{-j}t)\right)^{-1} \quad (t \in \mathbb{R}\setminus\{0\})
\]
即可. 由这个具体形式, 得 $\tau \in C^{1}(\mathbb{R})$. $\blacksquare$

\paragraph*{练习题 4.4}
(1) 设 $f \in L^{2}(\mathbb{R})$. 此时, $\|g*f\|_{L^{2}} \leq \|g\|_{L^{1}}\|f\|_{L^{2}}$ 成立. 因此 $\|L_{g}\|_{L^{2}\to L^{2}} \leq \|g\|_{L^{1}}$ 成立. 故 $L_{g}$ 有界.

(2) 对 $f, h \in L^{2}(\mathbb{R})$, 只需证明 $\langle L_{g}f, h\rangle = \langle f, L_{\check{g}}h\rangle$. 因 $L_{g}$, $L_{\check{g}}$ 有界, 由 $C_{\mathrm{c}}^{\infty}(\mathbb{R})$ 在 $L^{2}$ 中的稠密性, 可假定 $f, h \in C_{\mathrm{c}}^{\infty}(\mathbb{R})$. 此时, 由
\[
\int_{\mathbb{R}^{2}}|g(x-y)f(y)h(x)|\dd x\dd y \leq \|f\|_{L^{\infty}}\|h\|_{L^{1}}\|g\|_{L^{1}} < \infty
\]
可交换积分顺序, 对 $f, h \in C_{\mathrm{c}}^{\infty}(\mathbb{R})$, 得 $\langle L_{g}f, h\rangle = \langle f, L_{\check{g}}h\rangle$. 因此, $L_{g}$ 的伴随为 $L_{\check{g}}$. $\blacksquare$

\paragraph*{练习题 4.5}
由 $f, g \in C_{\mathrm{c}}^{1}(\mathbb{R})$, 存在某个 $R > 1$, 使得 $\operatorname{supp}(f) \subset [-R, R]$, $\operatorname{supp}(g) \subset [-R, R]$ 成立. 设 $F(t, s) = f(t-s)g(s)$, 由 $F_{t}(t, s) = f'(t-s)g(s)$, 得 $|F_{t}(t, s)| \leq \chi_{[-2R,2R]}(t)\max|f'|\max|g|$. 又 $|F(t, s)| \leq \chi_{[-2R,2R]}(t)\max|f|\max|g|$. 这里,
\begin{align*}
&\int_{\mathbb{R}}\chi_{[-2R,2R]}(t)\max|f'|\max|g|\dd t + \int_{\mathbb{R}}\chi_{[-2R,2R]}(t)\max|f|\max|g|\dd t \\
&= 2R(\max|f| + \max|f'|)\max|g| < \infty
\end{align*}
故由定理~\ref{thm:2.6}, 微分记号与积分记号可以交换, 得 $(f*g)' = f'*g$. 同理, 因 $f*g = g*f$, 也证明 $(f*g)' = f*g'$.

\paragraph*{练习题 4.6}
从略.

\paragraph*{练习题 4.7}
若 $I \in \mathcal{D}(\mathbb{R})$, 则 $I \subset [0, \infty)$ 或 $I \subset (-\infty, 0)$. 因此, 设 $x \in [0, \infty)$,
\begin{align*}
M_{\dyadic}f(x) &= \sup_{I \in \mathcal{D}(\mathbb{R})}\frac{\chi_{I}(x)}{|I|}\int_{I}|f(y)|\dd y \\
&= \sup_{I \in \mathcal{D}(\mathbb{R}),\, I \subset [0, \infty)}\frac{\chi_{I}(x)}{|I|}\int_{I}|f(y)|\dd y \\
&= \sup_{I \in \mathcal{D}(\mathbb{R})}\frac{\chi_{I}(x)}{|I|}\int_{I}\chi_{[0, \infty)}(y)|f(y)|\dd y \\
&= M_{\dyadic}[\chi_{[0, \infty)}f](x)
\end{align*}
即得. 同样, 设 $x \in (-\infty, 0)$,
\[
M_{\dyadic}f(x) = M_{\dyadic}[\chi_{(-\infty, 0)}f](x)
\]
即得. $\blacksquare$

\paragraph*{练习题 4.8}
因 $f$ 局部可积且 $Q$ 紧, 有 $\chi_{Q}f \in L^{1}(\mathbb{R}^{n})$, $\chi_{Q} \in L^{1}(\mathbb{R}^{n})$. 因此 $B$ 可积. 又
\[
\int_{Q}B(x)\dd x = \int_{Q}\chi_{Q}(x)f(x)\dd x - \frac{1}{|Q|}\int_{Q}f(z)\dd z \times \int_{Q}\chi_{Q}(x)\dd x = \int_{Q}\chi_{Q}(x)f(x)\dd x - \int_{Q}f(z)\dd z = 0.
\]
$\blacksquare$

\paragraph*{练习题 4.9}
(1) 证明逆否命题. 设 $\sup_{Q\in\mathcal{A}}|Q| = \infty$, 则对任意 $Q \in \mathcal{A}$, 因 $E \supset Q$, 有 $|E| \geq |Q|$, 由假定得 $|E| = \infty$.

(2) 设 $\mathcal{A}_{1} = \{Q \in \mathcal{A} : \text{存在 } R \in \mathcal{A} \text{ 使得 } R \ni Q\}$, 并定义 $\mathcal{A}_{0} = \mathcal{A}\setminus\mathcal{A}_{1}$. 于是, 一般地, 当 $Q_{1}, Q_{2} \in \mathcal{D}$ 不同时, $Q_{1}\cap Q_{2} = \emptyset$, $Q_{2} \subsetneq Q_{1}$, $Q_{1} \subsetneq Q_{2}$ 三者之一成立, 因此特别地, 当 $Q_{1}, Q_{2} \in \mathcal{A}_{0}$ 时, $Q_{1}\cap Q_{2} = \emptyset$. 又若 $Q \in \mathcal{A}_{1}$, 则存在 $R \in \mathcal{A}$ 使得 $Q \subset R$, 而由 (1), 这样的对象中存在体积最大者. 将其重新记为 $R$, 则 $R \in \mathcal{A}_{0}$. 因此,
\[
\bigcup_{Q \in \mathcal{A}}Q = \bigcup_{Q \in \mathcal{A}_{0}}Q \cup \bigcup_{Q \in \mathcal{A}_{1}}Q \subset \bigcup_{Q \in \mathcal{A}_{0}}Q \subset \bigcup_{Q \in \mathcal{A}}Q.
\]
即 $\bigcup_{Q\in\mathcal{A}}Q = \bigsqcup_{Q\in\mathcal{A}_{0}}Q$ 成立. 由此得结论.

\paragraph*{练习题 4.10}
当 $t > 2$ 时,
\[
H\chi_{[-1,1]}(t) = \int_{-1}^{1}\frac{1}{t-s}\dd s = [-\log(t-s)]_{-1}^{1} = \log\frac{1+t}{t-1}.
\]
由
\[
\lim_{t \to \infty}tH\chi_{[-1,1]}(t) = 2
\]
可知, 对充分大的 $t > 2$ (设 $t > t_{0} > 2$),
\[
|tH\chi_{[-1,1]}(t)| > 1.
\]
因此,
\[
\int_{t_{0}}^{\infty}|H\chi_{[-1,1]}(t)|\dd t \geq \int_{t_{0}}^{\infty}\frac{\dd t}{t} = \infty.
\]
故 $H\chi_{[-1,1]} \notin L^{1}([2,\infty))$. $\blacksquare$

\paragraph*{练习题 4.11}
当 $f \in L^{2}(\mathbb{R})$ 时, 由 $f_{j} = \chi_{B(j)}\chi_{[0,j]}(|f|)f$ 给出的函数满足 $|f - f_{j}|^{2} \leq |f|^{2}$. 因此, 由于 $\{f_{j}\}_{j=1}^{\infty}$ 逐点收敛到 $f$, 由 Lebesgue 收敛定理, $\{f_{j}\}_{j=1}^{\infty}$ 在 $L^{2}$ 中收敛到 $f$. 因 $f_{j} \in L^{p}(\mathbb{R}) \cap L^{2}(\mathbb{R})$, 故在 $L^{p}(\mathbb{R})$ 中稠密. $\blacksquare$

\paragraph*{练习题 4.12}
任意固定 $\varepsilon > 0$.
\[
\left|\int_{|y|>\varepsilon}\frac{f(x-y)+g(x-y)}{y}\dd y\right| \leq \left|\int_{|y|>\varepsilon}\frac{f(x-y)}{y}\dd y\right| + \left|\int_{|y|>\varepsilon}\frac{g(x-y)}{y}\dd y\right| \leq H^{*}f(x) + H^{*}g(x).
\]
因此, 对 $\varepsilon > 0$ 取上确界, 得 $H^{*}[f+g](x) \leq H^{*}f(x) + H^{*}g(x)$. $\blacksquare$

\paragraph*{练习题 4.13}
(1) 仔细考察定理~\ref{thm:4.6} 的证明.

(2) 因两边均为 0, 当 $T = 0$ 时结论显然. 假定并非如此. 此时, 因 $A+B+\|T\|_{L^{2}\to L^{2}} > 0$, 设
\[
S = \frac{T}{A+B+\|T\|_{L^{2}\to L^{2}}}
\]
则由 (1), 对每个 $f \in L^{p}(\mathbb{R}^{n})$, 得 $\|Sf\|_{L^{p}} \leq \zeta(1,1,1)\|f\|_{L^{p}}$. 因此, 对每个 $f \in L^{p}(\mathbb{R}^{n})$, 得 $\|Tf\|_{L^{p}} \leq \zeta(1,1,1)(A+B+\|T\|_{L^{2}\to L^{2}})\|f\|_{L^{p}}$. $\blacksquare$

\section*{第5章}

\paragraph*{练习题 5.1}
利用提示中的等式与 $\varphi'(t+k) \leq \varphi'(t)$, $t, k \geq 0$ 即可证明. $\blacksquare$

\paragraph*{练习题 5.2}
设 $\{f_{j}\}_{j=1}^{\infty}$ 是 $(\mathcal{S}(\mathbb{R}^{n}), d)$ 中的 Cauchy 列. 因 $\max_{x\in\mathbb{R}^{n}}|f_{j}(x)-f_{k}(x)| = p_{0,0}(f_{j}-f_{k})$, $\{f_{j}\}_{j=1}^{\infty}$ 一致收敛到某个连续函数 $f$. 同样, 对任意多重指数 $\alpha \in \mathbb{N}_{0}^{n}$, 由于 $\{\partial^{\alpha}f_{j}\}_{j=1}^{\infty}$ 是 $(\mathcal{S}(\mathbb{R}^{n}), d)$ 中的 Cauchy 列, $\{\partial^{\alpha}f_{j}\}_{j=1}^{\infty}$ 一致收敛. 因此 $\{\partial^{\alpha}f_{j}\}_{j=1}^{\infty}$ 一致收敛到 $\partial^{\alpha}f$, 从而 $\partial^{\alpha}f$ 连续. 特别地, $f \in C^{\infty}(\mathbb{R}^{n})$. 这里, 对任意多重指数 $\alpha, \beta \in \mathbb{N}_{0}^{n}$,
\[
p_{\alpha,\beta}(f-f_{j}) \leq \limsup_{k \to \infty}p_{\alpha,\beta}(f_{k}-f_{j})
\]
令 $j \to \infty$, 由定理~\ref{thm:5.1} 可知 $\{f_{j}\}_{j=1}^{\infty}$ 在 $(\mathcal{S}(\mathbb{R}^{n}), d)$ 中收敛. $\blacksquare$

\paragraph*{练习题 5.3}
(1) 由提示中的不等式,
\begin{align*}
\|f\|_{L^{p}}^{p} &= \int_{\mathbb{R}^{n}}\frac{1}{(1+|x|)^{p(n+1)}}\dd x \\
&\leq \int_{\mathbb{R}^{n}}\frac{1}{(1+|x|)^{n+1}}\dd x \\
&\leq \int_{\mathbb{R}^{n}}\prod_{j=1}^{n}\frac{1}{(1+|x_{j}|)^{\frac{n+1}{n}}}\dd x \\
&= \left(\int_{\mathbb{R}}\frac{1}{(1+|t|)^{\frac{n+1}{n}}}\dd t\right)^{n} \\
&= (2n)^{n} < \infty.
\end{align*}

(2) 由提示中的不等式, $M \leq q(g) + \sum_{j=1}^{n}r_{j}(g) < \infty$.

(3) 若 $g \in \mathcal{S}(\mathbb{R}^{n})$, 由 (2), $|g(x)| \leq \frac{M}{(1+|x|)^{n+1}}$. 由 (1), 右边属于 $L^{p}(\mathbb{R}^{n})$, 故 $g$ 也属于同一空间. $\blacksquare$

\paragraph*{练习题 5.4}
对函数 $\varphi$, 使用记号~\eqref{eq:5.1}. 取 $\alpha$, $\gamma$ 为多重指数, 由
\[
(f \cdot g)_{(\gamma, \alpha)} = \sum C_{\beta}f_{(\gamma, \beta)} \cdot g_{(0, \alpha-\beta)}
\]
得
\[
p_{\gamma,\beta}(f \cdot g) \leq \sum C_{\beta}p_{\gamma,\beta}(f)p_{0,\alpha-\beta}(g) < \infty.
\]
$\blacksquare$

\paragraph*{练习题 5.5}
由定理~\ref{thm:5.1}, 只需证明其与式~\eqref{eq:5.4} 对一切 $\alpha, \beta \in \mathbb{N}_{0}^{n}$ 满足
\[
\lim_{j \to \infty}p_{\alpha,\beta}(\varphi_{j}-\varphi) = 0
\]
是等价的. 若式~\eqref{eq:5.4} 成立, 则对一切 $\alpha, \beta \in \mathbb{N}_{0}^{n}$ 满足上式可由
\[
p_{\alpha,\beta}(\varphi_{j}-\varphi) \leq \sup_{x \in \mathbb{R}^{n}}(1+|x|)^{|\alpha|}|\partial^{\beta}\varphi(x) - \partial^{\beta}\varphi_{j}(x)|
\]
看出, 这是显然的. 来证明其逆. 对 $k_{1}, k_{2}, \ldots, k_{n} \in \mathbb{N}_{0}$, 记 $|k| = k_{1}+k_{2}+\cdots+k_{n}$, $x^{k} = x_{1}^{k_{1}}x_{2}^{k_{2}}\cdots x_{n}^{k_{n}}$, $k = (k_{1}, k_{2}, \ldots, k_{n})$, 由
\begin{align*}
(1+|x|)^{N} &\leq (1+|x_{1}|+|x_{2}|+\cdots+|x_{n}|)^{N} \\
&\leq \sum_{k \in \mathbb{N}_{0}^{n},\, |k| \leq N}\frac{(N-|k|)!k_{1}!k_{2}!\cdots k_{n}!}{N!}|x^{k}| \\
&\leq (N!)^{n}\sum_{k \in \mathbb{N}_{0}^{n},\, |k| \leq N}|x^{k}|
\end{align*}
得
\[
\sup_{x \in \mathbb{R}^{n}}(1+|x|)^{N}|\partial^{\alpha}\varphi(x) - \partial^{\beta}\varphi_{j}(x)| \leq (N!)^{n}\sum_{k \in \mathbb{N}_{0}^{n},\, |k| \leq N}\frac{(N-|k|)!k_{1}!k_{2}!\cdots k_{n}!}{N!}p_{k,\alpha}(\varphi_{j}-\varphi)
\]
故若 $\{\varphi_{j}\}_{j=1}^{\infty} \subset \mathcal{S}(\mathbb{R}^{n})$ 收敛到 $\varphi \in \mathcal{S}(\mathbb{R}^{n})$, 即对一切 $\alpha, \beta \in \mathbb{N}_{0}^{n}$ 满足上式, 则式~\eqref{eq:5.4} 成立. $\blacksquare$

\paragraph*{练习题 5.6}
$O$ 本身是开集族, 且 $\mathcal{S}'(\mathbb{R}^{n})$ 的任意开集族都包含在 $O$ 中, 故显然.

\paragraph*{练习题 5.7}
\begin{itemize}
\item $\mathcal{F}: \mathcal{S}'(\mathbb{R}^{n}) \to \mathcal{S}'(\mathbb{R}^{n})$ 是单射. 若 $\mathcal{F}f = 0$, 则 $\langle f, \varphi\rangle = \langle \mathcal{F}^{-1}\mathcal{F}f, \varphi\rangle = \langle \mathcal{F}f, \mathcal{F}^{-1}\varphi\rangle = \langle 0, \mathcal{F}^{-1}\varphi\rangle = 0$ 即得.

\item $\mathcal{F}: \mathcal{S}'(\mathbb{R}^{n}) \to \mathcal{S}'(\mathbb{R}^{n})$ 是满射. 任意给定 $F \in \mathcal{S}'(\mathbb{R}^{n})$, 定义 $G \in \mathcal{S}'(\mathbb{R}^{n})$ 为 $\langle G, \varphi\rangle = \langle F, \mathcal{F}^{-1}\varphi\rangle$, 则 $\langle \mathcal{F}G, \varphi\rangle = \langle G, \mathcal{F}\varphi\rangle = \langle F, \mathcal{F}^{-1}\mathcal{F}\varphi\rangle = \langle F, \varphi\rangle$, 故得 $F = \mathcal{F}G$.

\item $\mathcal{F}: \mathcal{S}'(\mathbb{R}^{n}) \to \mathcal{S}'(\mathbb{R}^{n})$ 连续. 只需证明在 $0 \in \mathcal{S}'(\mathbb{R}^{n})$ 处的连续性. 设 $\mathcal{Z} = \{g \in \mathcal{S}'(\mathbb{R}^{n}) : |\langle g, \varphi\rangle| < \varepsilon\}$. 这里, $\varepsilon > 0$, $\varphi \in \mathcal{S}(\mathbb{R}^{n})$. 于是, 因 $\mathcal{F}^{-1}[\mathcal{Z}] = \{g \in \mathcal{S}'(\mathbb{R}^{n}) : |\langle g, \mathcal{F}\varphi\rangle| < \varepsilon\}$, 可知 $\mathcal{F}^{-1}$ 连续.

\item $\mathcal{F}: \mathcal{S}'(\mathbb{R}^{n}) \to \mathcal{S}'(\mathbb{R}^{n})$ 的逆映射连续. 这可以与 $\mathcal{F}: \mathcal{S}'(\mathbb{R}^{n}) \to \mathcal{S}'(\mathbb{R}^{n})$ 连续同样地证明, 故省略. $\blacksquare$
\end{itemize}

\paragraph*{练习题 5.8}
在 $R = \pi$ 处取最大值. 因 $|f(t)| \leq f(\pi)$,
\[
\sup_{t, \varepsilon, R > 0}\left|\int_{\varepsilon}^{R}\frac{\sin(ts)}{s}\dd s\right| \leq 2f(\pi) < \infty.
\]
$\blacksquare$

\paragraph*{练习题 5.9}
设 $\varphi(x) = \frac{1}{(1+|x|)^{n+1}}$.

\begin{itemize}
\item 将例~\ref{exa:5.4} 应用于 $h(x) = \exp(-i\xi \cdot x)$, 结合 $\partial_{j}e^{-ix \cdot \xi} = -i\xi_{j}e^{-ix \cdot \xi}$, 可知式~\eqref{eq:5.7} 成立.

\item 设 $F(x, \xi) = f(x)e^{-ix \cdot \xi}$. 与练习题~\ref{exer:5.3}(2) 相同, $M = \sup_{x \in \mathbb{R}^{n}}(1+|x|)^{n+2}|f(x)| < \infty$. $|F(x, \xi)| \leq M(1+|x|)^{-n-2}$ 成立, 且由练习题~\ref{exer:5.3}(1), $\varphi \in L^{1}(\mathbb{R}^{n})$, 故
\[
\int_{\mathbb{R}^{n}}|F(x, \xi)|\dd x < \infty.
\]
又因 $F_{\xi_{j}}(x, \xi) = -ix_{j}f(x)e^{-ix \cdot \xi}$, $|F_{\xi_{j}}(x, \xi)| \leq M|x_{j}|(1+|x|)^{-n-2} \leq M\varphi(x)$ 成立, 再次由练习题~\ref{exer:5.3}(1), $\varphi \in L^{1}(\mathbb{R}^{n})$, 故可用定理~\ref{thm:2.13}, 式~\eqref{eq:5.8} 成立.

\item 由 $f \in \mathcal{S}(\mathbb{R}^{n})$, 式~\eqref{eq:5.7}, \eqref{eq:5.8} 可反复使用, 故利用绝对值为 1 的适当复常数 $c_{\alpha,\beta}$, 可作
\[
\xi^{\alpha}\partial^{\beta}\mathcal{F}\varphi(\xi) = c_{\alpha,\beta}\mathcal{F}(\partial^{\alpha}[x^{\beta}\varphi])(x)
\]
形式的变形. 因此只需考虑 $\alpha = \beta = 0$ 的情形. 当 $\alpha = \beta = 0$ 时, 由练习题~\ref{exer:5.3}(1), (2), 可取 $N = n+1$. $\blacksquare$
\end{itemize}

\paragraph*{练习题 5.10}
(1) 使用 $Ax \cdot \xi = x \cdot A^{\mathrm{t}}\xi = x \cdot A^{-1}\xi$ 与 $\det A = \pm 1$.

(2) 由于等式局部成立, 特别地, 取 $0 < \varepsilon < R < \infty$, 只需对 $\varepsilon < |x| < R$ 证明等式. 取满足 $0 \notin \operatorname{supp}(g)$ 且在 $\varepsilon < |x| < R$ 上取值为 1 的 $g \in C_{\mathrm{c}}^{\infty}(\mathbb{R}^{n})$, 考虑以乘积 $g \cdot u$ 代替 $u$. 因 $\|\nabla^{k}u\| = \|\nabla^{k}[g \cdot u]\|$, 若必要, 用 $g \cdot u$ 替换 $u$, 可假定 $u \in C_{\mathrm{c}}^{\infty}(\mathbb{R}^{n})$.
\[
\partial^{\alpha}u(x) = \partial_{\alpha_{1}}\partial_{\alpha_{2}}\cdots\partial_{\alpha_{k}}[\mathcal{F}^{-1}[\mathcal{F}u]](x) = \frac{(-i)^{k}}{(2\pi)^{\frac{n}{2}}}\int_{\mathbb{R}^{n}}e^{ix \cdot \xi}\xi^{\alpha}\mathcal{F}u(\xi)\dd\xi \quad (x \in \mathbb{R}^{n})
\]
以及
\[
\partial^{\alpha}u(x) = \partial_{\alpha_{1}}\partial_{\alpha_{2}}\cdots\partial_{\alpha_{k}}[\mathcal{F}[\mathcal{F}^{-1}u]](x) = \frac{i^{k}}{(2\pi)^{\frac{n}{2}}}\int_{\mathbb{R}^{n}}e^{-ix \cdot \eta}\eta^{\alpha}\mathcal{F}^{-1}u(\eta)\dd\eta \quad (x \in \mathbb{R}^{n}).
\]
因此,
\[
(2\pi)^{n}(\partial^{\alpha}u(x))^{2} = \int\int_{\mathbb{R}^{n} \times \mathbb{R}^{n}}e^{ix \cdot (\xi-\eta)}\xi^{\alpha} \cdot \eta^{\alpha}\mathcal{F}u(\xi)\mathcal{F}^{-1}u(\eta)\dd\xi\dd\eta.
\]
所以,
\begin{align*}
(2\pi)^{n}\|\nabla^{k}u(x)\|^{2} &= (2\pi)^{n}\sum_{\alpha \in \{1,2,\ldots,n\}^{k}}(\partial^{\alpha}u(x))^{2} \\
&= \sum_{\alpha \in \{1,2,\ldots,n\}^{k}}\int\int_{\mathbb{R}^{n} \times \mathbb{R}^{n}}e^{ix \cdot (\xi-\eta)}\xi^{\alpha} \cdot \eta^{\alpha}\mathcal{F}u(\xi)\mathcal{F}^{-1}u(\eta)\dd\xi\dd\eta \\
&= \int\int_{\mathbb{R}^{n} \times \mathbb{R}^{n}}e^{ix \cdot (\xi-\eta)}(\xi \cdot \eta)^{k}\mathcal{F}u(\xi)\mathcal{F}^{-1}u(\eta)\dd\xi\dd\eta \quad (x \in \mathbb{R}^{n}).
\end{align*}
这里, 设 $A \in O(n)$,
\[
Ax \cdot y = x \cdot {}^{\mathrm{t}}Ay, \quad Ax \cdot Ay = x \cdot y \quad (x, y \in \mathbb{R}^{n})
\]
成立. 由 (1),
\[
\mathcal{F}u(A\xi) = \mathcal{F}u(\xi), \quad \mathcal{F}^{-1}u(A\xi) = \mathcal{F}^{-1}u(\xi) \quad (\xi \in \mathbb{R}^{n}).
\]
由 $A$ 的正交性与变量替换,
\begin{align*}
\|\nabla^{k}u(Ax)\|^{2} &= \frac{1}{(2\pi)^{n}}\int\int_{\mathbb{R}^{n} \times \mathbb{R}^{n}}e^{iAx \cdot (\xi-\eta)}(\xi \cdot \eta)^{k}\mathcal{F}u(\xi)\mathcal{F}^{-1}u(\eta)\dd\xi\dd\eta \\
&= \frac{1}{(2\pi)^{n}}\int\int_{\mathbb{R}^{n} \times \mathbb{R}^{n}}e^{ix \cdot A(\xi-\eta)}(\xi \cdot \eta)^{k}\mathcal{F}u(\xi)\mathcal{F}^{-1}u(\eta)\dd\xi\dd\eta \\
&= \frac{1}{(2\pi)^{n}}\int\int_{\mathbb{R}^{n} \times \mathbb{R}^{n}}e^{ix \cdot (\tilde{\xi}-\tilde{\eta})}(A\tilde{\xi} \cdot A\tilde{\eta})^{k}\mathcal{F}u(A\tilde{\xi})\mathcal{F}^{-1}u(A\tilde{\eta})\dd\tilde{\xi}\dd\tilde{\eta} \\
&= \frac{1}{(2\pi)^{n}}\int\int_{\mathbb{R}^{n} \times \mathbb{R}^{n}}e^{ix \cdot (\tilde{\xi}-\tilde{\eta})}(\tilde{\xi} \cdot \tilde{\eta})^{k}\mathcal{F}u(\tilde{\xi})\mathcal{F}^{-1}u(\tilde{\eta})\dd\tilde{\xi}\dd\tilde{\eta} \\
&= \|\nabla^{k}u(x)\|^{2}.
\end{align*}
因此, $\|\nabla^{k}u\|^{2}$ 确实满足条件. $\blacksquare$

\paragraph*{练习题 5.11}
由分部积分
\begin{align*}
\left|\int_{\mathbb{R}}f(t)\varphi(t)\dd t\right| &= \left|\int_{\mathbb{R}}\varphi'(t)\cos(e^{t})\dd t\right| \\
&= \int_{\mathbb{R}}|\varphi'(t)\cos(e^{t})|\dd t \\
&\leq \pi p_{2,1}(\varphi)
\end{align*}
即得. 因此, 由 $C_{\mathrm{c}}^{\infty}(\mathbb{R}^{n})$ 在 $\mathcal{S}(\mathbb{R}^{n})$ 中的稠密性 (参照例~\ref{exa:5.7}) 得结论. $\blacksquare$

\paragraph*{练习题 5.12}
对任意 $\Phi \in \mathcal{S}(\mathbb{R}^{n})$, 只需证明
\[
\langle\mathcal{F}[f*g], \Phi\rangle = (2\pi)^{\frac{n}{2}}\langle\mathcal{F}g \cdot \mathcal{F}f, \Phi\rangle.
\]
由 $\mathcal{S}'(\mathbb{R}^{n})$ 中 Fourier 变换的定义,
\[
\langle\mathcal{F}[f*g], \Phi\rangle = \langle f*g, \mathcal{F}\Phi\rangle.
\]
这里, 设 $\Psi = \mathcal{F}\Phi(-\cdot)$, 由卷积的定义,
\[
\langle\mathcal{F}[f*g], \Phi\rangle = \langle f*g, \Psi(-\cdot)\rangle = (f*g)*\Psi(0).
\]
由定理~\ref{thm:5.12}(3),
\[
\langle\mathcal{F}[f*g], \Phi\rangle = g*(f*\Psi)(0) = \langle g, f*[\Psi(-\cdot)]\rangle.
\]
再次由 $\mathcal{S}(\mathbb{R}^{n})$ 中 Fourier 变换的定义,
\[
\langle\mathcal{F}[f*g], \Phi\rangle = \langle\mathcal{F}g, \mathcal{F}^{-1}[f*[\Psi(-\cdot)]]\rangle.
\]
这里,
\begin{align*}
\mathcal{F}^{-1}[f*[\Psi(-\cdot)]](\xi) &= \frac{1}{(2\pi)^{\frac{n}{2}}}\int_{\mathbb{R}^{n}}f*[\Psi(-\cdot)](x)e^{ix \cdot \xi}\dd x \\
&= \frac{1}{(2\pi)^{\frac{n}{2}}}\int_{\mathbb{R}^{n}}\left(\int_{\mathbb{R}^{n}}f(y)\Psi(-x+y)e^{ix \cdot \xi}\dd y\right)\dd x
\end{align*}
但由于 $f, \Psi \in \mathcal{S}(\mathbb{R}^{n}) \subset L^{1}(\mathbb{R}^{n})$, 可以交换积分顺序. 因此,
\[
\mathcal{F}[f*\Psi(-\cdot)](\xi) = \frac{1}{(2\pi)^{\frac{n}{2}}}\int_{\mathbb{R}^{n}}\left(\int_{\mathbb{R}^{n}}f(y)\Psi(-x+y)e^{ix \cdot \xi}\dd x\right)\dd y
\]
即得. 作变量替换,
\[
\mathcal{F}[f*\Psi(-\cdot)](\xi) = (2\pi)^{\frac{n}{2}}\mathcal{F}f(\xi)\mathcal{F}\Psi(\xi) = (2\pi)^{\frac{n}{2}}\mathcal{F}f(\xi)\Phi(\xi)
\]
将其代入, 运用 $\mathcal{F}$ 的定义, 得
\[
\langle\mathcal{F}[f*g], \Phi\rangle = (2\pi)^{\frac{n}{2}}\langle\mathcal{F}g \cdot \mathcal{F}f, \Phi\rangle.
\]
$\blacksquare$

\section*{第6章}

\paragraph*{练习题 6.1}
仿照引理~\ref{lem:6.2} 的证明即可. $\blacksquare$

\paragraph*{练习题 6.2}
因三个线性算子 $\mathcal{F}$, $f \mapsto mf$, $\mathcal{F}^{-1}$ 都是 $L^{2}$ 有界的, 故 $T_{m}$ 是 $L^{2}$ 有界的. $\blacksquare$

\paragraph*{练习题 6.3}
$\mathcal{F}[T_{m}g] = m\mathcal{F}g$. 对任意 $\alpha, \beta \in \mathbb{N}_{0}^{n}$ 与 $N \in \mathbb{N}$, $\partial^{\alpha}m(\xi) = O(|\xi|^{-|\alpha|})$ 且 $\partial^{\beta}\mathcal{F}g(\xi) = O(|\xi|^{N})$, 故由 Leibniz 法则, $\partial^{\alpha}\mathcal{F}[T_{m}g] = O(|\xi|^{N})$ 成立. 因此 $T_{m}g \in \mathcal{S}_{\infty}$. $\blacksquare$
